\documentclass[10pt,a4paper]{amsart}
\usepackage[margin=19mm]{geometry}
\usepackage{amsmath,amssymb,mathtools}
\usepackage[scr=rsfs,cal=euler]{mathalfa}
\usepackage{xfrac}
\usepackage[unicode,hidelinks,bookmarks=false]{hyperref}
\newcommand{\ie}{i.e.\ }
\newcommand{\cf}{cf.\ }
\newcommand{\resp}{respectively\ }
\newcommand{\Subsection}{\subsection}
\providecommand{\nobreakdash}{\nobreak\hbox{-}}
\setlength{\emergencystretch}{2em}
\multlinegap0pt
\usepackage{mathrsfs,amscd,enumitem,amscd,color}
\usepackage{tikz}
\def\thefootnote{}
\usepackage{lmodern}
\newtheorem{ThA}{Theorem}
\renewcommand{\theThA}{\Alph{ThA}}
\newtheorem{thm}{Theorem}[section]
 \newtheorem{cor}[thm]{Corollary}
 \newtheorem{lem}[thm]{Lemma}
 \newtheorem{prop}[thm]{Proposition}
\newtheorem{defn}[thm]{Definition}
 \newtheorem{rem}[thm]{Remark}
\newtheorem{af}[thm]{Claim}
\newcommand{\im}{\textrm{Im}}
\newcommand{\N}{\mathbb{N}}
\newcommand{\Z}{\mathbb{Z}}
\newcommand{\LL}{\mathcal{L}}
\newcommand{\Cn}{\mathbb{C}^n}
\newcommand{\Rm}{\mathbb{R}_+}
\newcommand{\R}{\mathbb{R}}
\newcommand{\Rn}{\mathbb{R}^n}
\newcommand{\supp}{\mathop{\mathrm{supp}}}
\newcommand{\card}{\mathop{\mathrm{card}}}
\newcommand{\diam}{\mathop{\mathrm{diam}}}
\newcommand{\esssup}{\mathop{\mathrm{ess\, sup \;}}}
\newcommand{\essinf}{\mathop{\mathrm{ess\, inf \;}}}
\numberwithin{equation}{section}
\allowdisplaybreaks
\begin{document}
\title[Higher order logarithms of Bessel operators and an extension]{Higher order logarithms of Bessel operators and an extension problem}
\author{Jorge J. Betancor, Estefanía. Dalmasso, Juan C. Fariña; Pablo Quijano}
\date{}
\maketitle
\noindent\emph{Source and licence.} Higher order logarithms of Bessel operators and an extension problem. Version arXiv:2608.19516v1 (CC BY 4.0). This material is distributed under the Creative Commons Attribution 4.0 International licence, http://creativecommons.org/licenses/by/4.0/, as recorded in the arXiv OAI licence metadata for this exact version and shown on the abstract page https://arxiv.org/abs/2608.19516v1. Full rights notice: licenses/arxiv-2608.19516-CC-BY-4.0.txt.\par\smallskip
\noindent\textit{Editorial scope.} Counted unit: the original \texttt{prop} environment \texttt{prop: Prop2.5} at source lines 561--581 together with its complete proof at lines 583--613; the delivered document keeps source lines 426--614 so that every referenced label and every citation resolves inside it. Native editable source is the paper's own concatenated TeX; the AMS wrapper is editorial formatting only. No publisher class, upstream script or unrelated artwork is redistributed.
\section{Fractional powers of Bessel operators}\label{sec: fract powers}

In this section we establish representations through the semigroup $\{W_t^\lambda \}_{t>0}$ for the fractional powers $B_\lambda^s$, $s\in (-1,1)$, $s\neq 0$, of $B_\lambda$.


Let $s>0$. We define the operators $\mathbb B_{\lambda ,s}$ and $\mathbb B_{\lambda ,-s}$ as follows:
\[
\mathbb B_{\lambda, s}f=\frac{1}{\Gamma (-s)}\int_0^\infty \frac{W_t^\lambda (f)-f}{t^{1+s}}dt,\quad f\in D(\mathbb B_{\lambda, s}),
\]
where
\[
D(\mathbb B_{\lambda, s})=\left\{f\in L^2(0,\infty): \int_0^\infty \frac{\|W_t^\lambda (f)-f\|_{L^2(0,\infty)}}{t^{1+s}}dt<\infty\right\},
\]
and 
\[
\mathbb B_{\lambda, -s}f=\frac{1}{\Gamma (s)}\int_0^\infty W_t^\lambda (f)t^{s-1}dt,\quad f\in D(\mathbb B_{\lambda, -s}),
\]
being
\[
D(\mathbb B_{\lambda, -s})=\left\{f\in L^2(0,\infty): \int_0^\infty \|W_t^\lambda (f)\|_{L^2(0,\infty)}t^{s-1}dt<\infty\right\}.
\]
Note that the function $F:[0,\infty)\to L^2(0,\infty)$ given by $F(u)=W_u^\lambda (f)$, with $f\in L^2(0,\infty)$, is continuous, because $\{W_t^\lambda \}_{t>0}$ is a $C_0$-semigroup in $L^2(0,\infty)$. Since $L^2(0,\infty)$ is a separable Banach space, according to Pettis' Theorem (\cite[p. 131]{Yo}), the function $F$ is strongly measurable. Hence, if $f\in D(\mathbb B_{\lambda ,s})$ (respectively, $D(\mathbb B_{\lambda ,-s})$) the function $t\in (0,\infty)\mapsto (W_t^\lambda (f)-f)/t\in L^2(0,\infty)$ (respectively, $t\in (0,\infty)\mapsto W_t^\lambda (f)t^{s-1}\in L^2(0,\infty)$) is $L^2(0,\infty)$-Bochner integrable in $(0,\infty)$. According to \cite[(3), p. 260]{Yo}, $B_\lambda^sf=\mathbb B_{\lambda ,s}f$, $f\in D(B_\lambda)$. 

We now establish relations between $B_\lambda^s$ and $B_\lambda^{-s}$ with $\mathbb B_{\lambda , s}$ and $\mathbb B_{\lambda , -s}$, respectively.


\begin{prop}\label{prop: Prop2.1}
  Let $\lambda >-1$ and $s_0\in (0,1)$. If $f\in D(B_\lambda^{s_0})$, then $f\in D(\mathbb B_{\lambda ,s})$ and $\mathbb B_{\lambda ,s}f=B_\lambda^sf$, $s\in (0,s_0)$.  
\end{prop}
\begin{proof}
Let $s\in (0,s_0)$. Suppose that $f\in D(B_\lambda^{s_0})$. We have that $f\in L^2(0,\infty)$ and $y^{2{s_0}}h_\lambda (f)\in L^2(0,\infty)$. Then, $f\in D(B_\lambda^s)$. Since $h_\lambda $ is an isometry in $L^2(0,\infty)$ we get
\[
\|W_t^\lambda (f)-f\|_{L^2 (0,\infty)}=\|h_\lambda (W_t^\lambda (f))-h_\lambda (f)\|_{L^2(0,\infty)}=\|(e^{-ty^2}-1)h_\lambda (f)\|_{L^2(0,\infty)},\quad t>0.
\]
Then,
\begin{align*}
    \int_0^\infty \frac{\|W_t^\lambda (f)-f\|_{L^2(0,\infty)}}{t^{1+s}}dt&\leq C\left(\int_0^1\frac{\|(ty^2)^{s_0}h_\lambda (f)\|_{L^2(0,\infty)}}{t^{1+s}}dt+\int_1^\infty \frac{\|h_\lambda (f)\|_{L^2(0,\infty)}}{t^{1+s}}dt\right)\\
    &\leq C\left(\|y^{2s_0}h_\lambda (f)\|_{L^2(0,\infty)}+\|f\|_{L^2(0,\infty)}\right).
\end{align*}
It follows that $f\in D(\mathbb B_{\lambda ,s})$, and using properties of the Bochner integrals we get
\begin{align*}
\mathbb B_{\lambda ,s}f&=\frac{1}{\Gamma (-s)}\int_0^\infty \frac{W_t^\lambda (f)-f}{t^{1+s}}dt=\frac{1}{\Gamma (-s)}\int_0^\infty\frac{h_\lambda (e^{-y^2t}h_\lambda (f))-h_\lambda (h_\lambda (f))}{t^{1+s}}dt\\
    &=h_\lambda \left(\frac{1}{\Gamma (-s)}\int_0^\infty \frac{e^{-ty^2}-1}{t^{1+s}}dt\,h_\lambda (f)\right)=h_\lambda (y^{2s}h_\lambda (f))=B_\lambda^sf,
\end{align*}
because $\int_0^\infty \frac{e^{-zt}-1}{t^{1+s}}dt=z^s\Gamma (-s)$, $z>0$.
\end{proof}

\begin{prop}\label{prop: Prop2.2}
 Let $\lambda >-1$ and $s_0\in (0,1)$. If $f\in D(B_\lambda^{-s_0})$, then $f\in D(\mathbb B_{\lambda ,-s})$ and $\mathbb B_{\lambda ,-s}f=B_\lambda^{-s}f$, $s\in (0,s_0)$.     
\end{prop}
\begin{proof}
    Let $s\in (0,s_0)$. Assume that $f\in D(B_\lambda^{-s_0})$. Then, $f\in D(B_\lambda^{-s})$. We get
\[
\|W_t^\lambda (f)\|_{L^2 (0,\infty)}=\|h_\lambda (e^{-ty^2}h_\lambda (f))\|_{L^2(0,\infty)}=\|e^{-ty^2}h_\lambda (f)\|_{L^2(0,\infty)},\quad t>0.
\]
It follows that
\begin{align*}
    \int_0^\infty \frac{\|W_t^\lambda (f)\|_{L^2(0,\infty)}}{t^{1-s}}dt&=\int_0^\infty \frac{\|e^{-ty^2}h_\lambda (f)\|_{L^2(0,\infty)}}{t^{1-s}}dt\\
    &\leq \|f\|_{L^2(0,\infty)}\int_0^1\frac{dt}{t^{1-s}}+C\int_1^\infty \frac{\|y^{-2s_0}h_\lambda (f)\|_{L^2(0,\infty)}}{t^{1-s+s_0}}dt\\
    &\leq C\Big(\|f\|_{L^2(0,\infty)}+\|y^{-2s_0}h_\lambda (f)\|_{L^2(0,\infty)}\Big).
\end{align*}
Then, $f\in D(\mathbb B_{\lambda ,-s})$ and 
\[
\mathbb B_{\lambda ,-s}f=\frac{1}{\Gamma (s)}\int_0^\infty \frac{W_t^\lambda (f)}{t^{1-s}}dt=h_\lambda \left(\frac{1}{\Gamma (s)}\int_0^\infty e^{-ty^2}t^{s-1}dt\,h_\lambda (f)\right)=h_\lambda (y^{-2s}h_\lambda (f))=B_\lambda^{-s}f.\qedhere
\]
\end{proof}
We now study the derivatives $\partial_s^m B_\lambda^s$ of $B_\lambda^s$ in $L^2(0,\infty)$ with $m\in \mathbb N$ and $s\in (-1,1)$, $s\neq 0$.

\begin{prop}\label{prop: Prop2.3}
    Let $\lambda >-1$, $s_0\in (0,1)$ and $m_0\in \mathbb N$. Suppose that $f\in D(B_\lambda^{s_0})\cap D(\log^{m_0}B_\lambda)$. Then
    \[
    \partial_s^m B_\lambda^sf=(\log^m B_\lambda)B_\lambda^sf,\quad s\in (0,s_0)\mbox{ and }m\in \mathbb N,\,m\leq m_0.
    \] 
 Here, $\partial_s$ is understood in $L^2(0,\infty)$.
\end{prop}

\begin{proof}
    Let $s\in (0,s_0)$ and let $m\in \mathbb N$, $m \leq m_0$. Since $D( B_\lambda^{s_0})\subseteq D(B_\lambda^s)$, $f\in D(B_\lambda^s)$ and $B_\lambda^sf=h_\lambda (y^{2s}h_\lambda (f))$. Moreover, we have that
    \begin{align*}
        |\log y|^m y^{2s}&\leq C\left(|\log y|^m\mathcal X_{(0,\tfrac12)}(y)+y ^{2s_0}\mathcal X_{(\tfrac12,\infty)}(y)\right)\\
        &\leq C\left(|\log y|^{m_0}+y^{2s_0}\right),\quad y\in (0,\infty).
    \end{align*}
Then, $B_\lambda^sf\in D(\log^m B_\lambda)$.

By using the mean value theorem we get
\[
|y^{2h}-1|\leq C\left(1+y^{2(s_0-s)}\right)h\log y,\quad y\in (0,\infty)\mbox{ and }h\in (0,s_0-s).
\]
Since $B_\lambda^sf\in D(\log B_\lambda)$ and $B_\lambda^{s_0-s}f\in D(\log B_\lambda)$, by using the dominated convergence theorem we obtain
\[
\lim_{h\rightarrow 0^+}\int_0^\infty \left|\left(\frac{y^{2h}-1}{h}-2\log y\right)y^{2s}h_\lambda (f)(y)\right|^2dy=0.
\]
We also have that
\[
|y^{-2h}-1|\leq 2h(1+y^{-s})\log y,\quad y\in (0,\infty)\mbox{ and }h\in \big(0,\tfrac{s}{2}\big),
\]
and
\[
\lim_{h\rightarrow 0^-}\int_0^\infty \left|\left(\frac{y^{2h}-1}{h}-2\log y\right)y^{2s}h_\lambda (f)(y)\right|^2dy=0.
\]
Thus, we proved that
\[
\lim_{h\rightarrow 0}\frac{B_\lambda^{s+h}f-B_\lambda^sf}{h}=(\log B_\lambda)B_\lambda^sf,\quad \mbox{ in }L^2(0,\infty).
\]
By proceeding in a similar way we can see that if $m\in \mathbb N$, $m\leq m_0-1$, and $\partial_s^m B_\lambda^sf=(\log^m B_\lambda)B_\lambda^sf$, $0<s<s_0$, then $\partial_s^{m+1}B_\lambda^sf=(\log^{m+1}B_\lambda)B_\lambda^s f$, $0<s<s_0$.

Thus the proof is finished.
\end{proof}
By arguing as in the proof of the previous proposition we obtain the following result.

\begin{prop}\label{prop: Prop2.4}
    Let $\lambda >-1$, $s_0\in (0,1)$ and $m_0\in \mathbb N$. Suppose that $f\in D(B_\lambda^{-s_0})\cap D(\log^{m_0}B_\lambda)$. Then, 
    \[
    \partial_s^m B_\lambda^{-s}f=(-1)^m(\log^m B_\lambda)B_\lambda^{-s}f,\quad s\in (0,s_0)\mbox{ and }m\in \mathbb N,\,m\leq m_0,
    \] 
    where $\partial_s$ is understood in $L^2(0,\infty)$.
\end{prop}

From Propositions \ref{prop: Prop2.3} and \ref{prop: Prop2.4} we can deduce the following property that can be established using the dominated convergence theorem.

\begin{cor}\label{cor: Cor2.5}
    Let $\lambda >-1$, $s_0\in (0,1)$, and $m_0\in \mathbb N$.
\begin{enumerate} 
    \item If $f\in D(B_\lambda^{s_0})\cap D(\log^{m_0}B_\lambda)$ then for every $m\in \mathbb N$, $m\leq m_0$,
\[
(\log^m B_\lambda)f=\lim_{s\rightarrow 0^+}\partial_s^m B_\lambda^sf,\quad \mbox{ in }L^2(0,\infty).
\]
\item If $f\in D(B_\lambda^{-s_0})\cap D(\log^{m_0}B_\lambda)$ then for every $m\in \mathbb N$, $m\leq m_0$,
\[
(\log^m B_\lambda)f=\lim_{s\rightarrow 0^+}(-1)^m\partial_s^m B_\lambda^{-s}f,\quad \mbox{ in }L^2(0,\infty).
\]
\end{enumerate}
\end{cor}
We now establish asymptotic expansions for $B_\lambda^s$, $s\in (-1,1)$, $s\neq 0$, through logarithms of $B_\lambda$ in the $L^2(0,\infty)$-setting.


\begin{prop}\label{prop: Prop2.5}
  Let $\lambda >-1$, $s_0\in (0,1)$, $0<s<s_0$, and $m_0\in \mathbb N$.
  \begin{enumerate}
      \item \label{itm: Prop2.5 a} Suppose that $f\in D(B_\lambda^{s_0})\cap D(\log^{m_0+1}B_\lambda)$. Then,
      \[
      B_\lambda^s f=f+\sum_{j=1}^{m_0}\frac{s^j}{j!}(\log^jB_\lambda)f+o(s^{m_0}),\quad \mbox{ in }L^2(0,\infty),
      \]
      and
      \[
      \lim_{s\rightarrow 0^+}\frac{(m_0+1)!}{s^{m_0+1}}\left(B_\lambda^sf-f-\sum_{j=1}^{m_0}\frac{s^j}{j!}(\log^jB_\lambda)f\right)=(\log^{m_0+1}B_\lambda)f,\quad \mbox{ in }L^2(0,\infty).
      \]
      \item \label{itm: Prop2.5 b} Suppose that $f\in D(B_\lambda^{-s_0})\cap D(\log^{m_0+1}B_\lambda)$. Then,
      \[
      B_\lambda^{-s}f=f+\sum_{j=1}^{m_0}\frac{(-1)^js^j}{j!}(\log^jB_\lambda)f+o(s^{m_0}),\quad \mbox{ in }L^2(0,\infty),
      \]
      and
      \[
      \lim_{s\rightarrow 0^+}\frac{(m_0+1)!}{s^{m_0+1}}\left(B_\lambda^{-s}f-f-\sum_{j=1}^{m_0}\frac{(-1)^js^j}{j!}(\log^jB_\lambda)f\right)=(-1)^{m_0+1}(\log^{m_0+1}B_\lambda)f,\quad \mbox{ in }L^2(0,\infty).
      \]
  \end{enumerate}
\end{prop}

\begin{proof}
    We prove \ref{itm: Prop2.5 a}. We have
    \[
     B_\lambda^sf-f-\sum_{j=1}^{m_0}\frac{s^j}{j!}(\log^jB_\lambda)f=h_\lambda \left(\left(y^{2s}
     -1-\sum_{j=1}^{m_0}\frac{s^j}{j!}(2\log y)^j\right)h_\lambda (f)\right),\quad s\in (0,s_0).
    \]
    Since $h_\lambda$ is an isometry on $L^2(0,\infty)$ it follows that
    \[
    \left\|B_\lambda^sf-f-\sum_{j=1}^{m_0}\frac{s^j}{j!}(\log^jB_\lambda)f\right\|_{L^2(0,\infty)}=\left\|\left(y^{2s}
     -1-\sum_{j=1}^{m_0}\frac{s^j}{j!}(2\log y)^j\right)h_\lambda (f)\right\|_{L^2(0,\infty)},\quad s\in (0,s_0).
    \]
 Using the mean value theorem we obtain, for every $s\in (0,s_0)$ and $y\in (0,\infty)$,
    \[
    y^{2s}-1-\sum_{j=1}^{m_0}\frac{s^j}{j!}(2\log y)^j=(2\log y)^{m_0+1}y^{2u}\frac{s^{m_0+1}}{(m_0+1)!},
    \]
     for a certain $u\in (0,s)$. Then,
     \[
     \left|y^{2s}-1-\sum_{j=1}^{m_0}\frac{s^j}{j!}(2\log y)^j\right|\leq C\left(|\log y|^{m_0+1}+y^{2s_0}\right)s^{m_0+1},\quad s\in (0,s_0)\mbox{ and }y\in (0,\infty).
     \]
     We get, for each $s\in (0,s_0)$,
     \[
     \left\|B_\lambda^sf-f-\sum_{j=1}^{m_0}\frac{s^j}{j!}(\log^jB_\lambda)f\right\|_{L^2(0,\infty)}\leq C\left(\|B_\lambda^{s_0}f\|_{L^2(0,\infty)}+\|(\log^{m_0+1}B_\lambda)f\|_{L^2(0,\infty)}\right)s^{m_0+1}.
     \]
     The dominated convergence theorem leads to 
     \begin{align*}
         \lim_{s\rightarrow 0^+}\left\|\frac{(m_0+1)!}{s^{m_0+1}}\left(B_\lambda^sf-f-\sum_{j=1}^{m_0}\frac{s^j}{j!}(\log^jB_\lambda)f\right)-(\log^{m_0+1}B_\lambda)f\right\|_{L^2(0,\infty)}\\
         &\hspace{-10cm}=\lim_{s\rightarrow 0^+}\left\|\left[\frac{(m_0+1)!}{s^{m_0+1}}\left(y^{2s}-1-\sum_{j=1}^{m_0}\frac{(2s)^j}{j!}(\log y)^j\right)-(2\log y)^{m_0+1}\right]h_\lambda (f)\right\|_{L^2(0,\infty)}=0.
     \end{align*}

     The property in \ref{itm: Prop2.5 b} can be established in a similar way.
\end{proof}


\begin{thebibliography}{99}
\bibitem[Yo]{Yo} {\sc Yosida, K.} \newblock {\em Functional analysis}, sixth~ed., vol.~123 of {\em Grundlehren der Mathematischen Wissenschaften}. \newblock Springer-Verlag, Berlin-New York, 1980.
\end{thebibliography}
\end{document}
